% Complete finite proofs for the first structural causal model component.
\section{Finite structural causal models: proofs}
\label{sec:annex-p16-finite-causal-models}

Use the graph, probability and equation-replacement definitions in
\cref{sec:p16-finite-causal-models}. The source supplies the structural
interpretation, intervention operation, identification criterion and
factorization statements
\cite[secs.~3.1--3.2.3, pp.~103--112]{Pearl2009CausalOverview}.
The finite constructions and proofs here include explicit treatment of
zero probabilities and dependent disturbances.

\begin{annexitem}{annex:p16.causal-topological-order}{thm:p16-causal-topological-order}{Ordering an acyclic finite graph}
\begin{proof}
First consider a nonempty finite acyclic directed graph. Suppose that every
vertex had an incoming edge. Choose a vertex \(j_0\), then choose an
incoming predecessor \(j_1\to j_0\), then a predecessor
\(j_2\to j_1\), and continue until one more vertex has been selected
than the graph contains. Some vertex repeats, say \(j_a=j_b\)
with \(a<b\). The edges
\(j_b\to j_{b-1}\to\cdots\to j_a\) form a directed cycle,
contrary to acyclicity. Therefore some vertex has no incoming edge.

Remove such a vertex and all its incident edges. The remaining graph is
acyclic, because any cycle in it would be a cycle in the original graph.
Whenever the remaining graph is nonempty, the preceding argument supplies
another vertex to remove. Exactly one distinct vertex is removed at each
step, so the procedure stops after \(n\) steps. Let
\(i_1,\ldots,i_n\) be the removal order. If \(j\to i_r\), then
\(j\) cannot still be in the graph when \(i_r\) is removed, because
it would supply an incoming edge. Thus \(j\) occurs earlier in the
order. This proves the topological property and termination.
\end{proof}
\end{annexitem}

\begin{annexitem}{annex:p16.causal-solution}{thm:p16-causal-solution}{Unique solutions and probability laws}
\begin{proof}
Fix \(\omega\in\Omega\) and a topological order
\(i_1,\ldots,i_n\). At the first vertex the parent tuple is empty.
The total function \(f_{i_1}\), applied to that tuple and
\(U_{i_1}(\omega)\), gives a member of \(S_{i_1}\). At step
\(r\), every parent of \(i_r\) has already been assigned a value,
so the total function \(f_{i_r}\) gives its value. Subsequent steps
never change earlier coordinates. Consequently every structural equation
is satisfied by the completed tuple.

If \(v,w\) are two solutions at \(\omega\), their first
coordinates in the topological order agree, since each applies the same
function to the same empty parent tuple and the same disturbance. If their
earlier coordinates agree, the parents of \(i_r\) agree and the
equation for \(i_r\) forces its values to agree. Induction gives
\(v=w\). Any other topological evaluation is also a solution, hence
must give this same tuple. Since this holds separately at every outcome,
it defines functions \(V_i:\Omega\to S_i\).

For each \(v\in S_N\), the set \(\{\omega:V(\omega)=v\}\)
is an event in the finite probability space. These sets are pairwise
disjoint and their union is \(\Omega\); empty fibers contribute zero.
Thus each stated mass is nonnegative and
\[
 \sum_{v\in S_N}p_M(v)
 =\sum_{v\in S_N}\sum_{\omega:V(\omega)=v}q(\omega)
 =\sum_{\omega\in\Omega}q(\omega)=1.
\]
For a fixed \(J\subseteq N\) and \(s\in S_J\), partition the
event \(\{V_J=s\}\) by its full assignments \(v\) with
\(v_J=s\). Finite additivity gives the marginal formula. This
argument also covers \(J=\varnothing\): there is just one empty
tuple, of probability one. No division by a probability is involved.
\end{proof}
\end{annexitem}

\begin{annexitem}{annex:p16.causal-intervention-laws}{thm:p16-causal-intervention-laws}{Equation replacement, consistency and unaffected vertices}
\begin{proof}
Deleting incoming edges cannot introduce a directed cycle, so the
intervention graph is again a DAG. The replacement functions are total
constant functions into their original value sets. All other functions
and the finite probability space remain valid model data. The preceding
solution theorem therefore applies to \(M_{I=x}\), proving its
unique solution and the mass formula. The equations at vertices in
\(I\) force its solution to have coordinates \(x\), so assignments
with different coordinates have empty fibers and probability zero.

For the first pointwise assertion, suppose \(V_I(\omega)=x\).
The original tuple \(V(\omega)\) already satisfies every unchanged
equation, and by the supposition it also satisfies every replacement
equation. It is therefore a solution of the modified system at
\(\omega\). Uniqueness implies
\(V^{I=x}(\omega)=V(\omega)\).

For the second assertion, put
\(H=N\setminus(I\cup\operatorname{de}_G(I))\). Every parent of
a vertex in \(H\) is in \(H\): a parent in \(I\) gives a
length-one path from \(I\), and a parent in
\(\operatorname{de}_G(I)\) gives such a path after appending its
edge. Either possibility would contradict membership of the child in
\(H\). Read the original topological order restricted to \(H\).
At its first vertex the parent inputs, if any, cannot occur outside
\(H\) or later in the order. At each vertex in \(H\), its equation
and disturbance are unchanged, and its parent values agree by induction.
The value at that vertex therefore agrees in the two evaluations. This
proves the pointwise assertion throughout \(H\).

The empty intervention replaces no equation and deletes no edge. For
disjoint \(I,J\), each original equation belongs to exactly one of
three cases: it is replaced by the prescribed constant in \(I\),
replaced by the prescribed constant in \(J\), or left unchanged.
The order of these two replacements does not change any case, graph,
probability weight or disturbance function. The resulting models are the
same, so their solutions and distributions are the same as well.
\end{proof}
\end{annexitem}

\begin{annexitem}{annex:p16.causal-factorization}{thm:p16-causal-factorization}{Independent disturbances give the structural product}
\begin{proof}
We first record the finite consequence of the stated mutual independence.
For arbitrary subsets \(A_i\subseteq\mathcal U_i\), partition the
event \(\bigcap_i\{U_i\in A_i\}\) by its disturbance tuples.
Then finite distributivity gives
\[
\begin{aligned}
 \mathbb P\Bigl(\bigcap_{i\in N}\{U_i\in A_i\}\Bigr)
 &=\sum_{u_1\in A_1}\cdots\sum_{u_n\in A_n}
                         \prod_{i\in N}r_i(u_i)\\
 &=\prod_{i\in N}\sum_{u_i\in A_i}r_i(u_i).
\end{aligned}
\]
An empty set gives zero on both sides. Taking \(A_i=\mathcal U_i\)
for omitted indices proves the same identity for a subset of the
disturbances, because those factors sum to one.

Fix \(v\in S_N\), and define
\[
 A_i(v)=\{u\in\mathcal U_i:
               f_i(v_{\operatorname{pa}_G(i)},u)=v_i\}.
\]
For each \(\omega\), the assertion \(U_i(\omega)\in A_i(v)\)
for every \(i\) says exactly that the fixed tuple \(v\) satisfies
all structural equations at that outcome. By uniqueness this holds
exactly when \(V(\omega)=v\). Therefore
\[
 p_M(v)=\mathbb P\Bigl(\bigcap_i\{U_i\in A_i(v)\}\Bigr)
       =\prod_i\sum_{u\in A_i(v)}r_i(u)
       =\prod_i k_i(v_i\mid v_{\operatorname{pa}_G(i)}).
\]
This proves the product at every tuple, without conditioning.

For the conditional assertion, take a topological order in which vertex
\(i\) occurs in position \(r\), and let \(L\) be the set of
vertices in earlier positions. Recursive evaluation shows that the tuple
\(V_L\) is a function of \((U_j)_{j\in L}\): at each earlier
step its parents also lie earlier, so substituting their already
constructed functions introduces no disturbance outside \(L\).
In particular \(V_{\operatorname{pa}_G(i)}\) is a function of
this disturbance tuple, and \(i\notin L\).

For a tuple \(u_L\), summing the joint disturbance mass over all
indices outside \(L\cup\{i\}\) gives
\[
 \mathbb P(U_i=u,(U_j)_{j\in L}=u_L)
       =r_i(u)\prod_{j\in L}r_j((u_L)_j).
\]
When \(L\) is empty, its tuple has probability one and the same
identity holds. Sum this identity over the tuples \(u_L\) whose
evaluation gives a specified parent assignment \(a\). These tuple
events partition \(\{V_{\operatorname{pa}_G(i)}=a\}\). Hence
\[
 \mathbb P(U_i=u,V_{\operatorname{pa}_G(i)}=a)
       =r_i(u)\mathbb P(V_{\operatorname{pa}_G(i)}=a).
\]
On this parent event, the equation for \(i\) gives
\(V_i=f_i(a,U_i)\). Summing the last identity over the disturbance
values with output \(b\) yields
\[
 \mathbb P(V_i=b,V_{\operatorname{pa}_G(i)}=a)
       =k_i(b\mid a)\mathbb P(V_{\operatorname{pa}_G(i)}=a).
\]
Division by the assumed positive final factor proves the conditional
identity. If that factor is zero, the displayed joint identity still
holds, but it supplies no conditional ratio. The structural kernel
remains defined by its original finite sum.
\end{proof}
\end{annexitem}

\begin{annexitem}{annex:p16.causal-truncation}{thm:p16-causal-truncation}{Truncating the product after replacement}
\begin{proof}
The modified model has the same disturbance functions and their same
joint law, so mutual independence still holds. For \(i\notin I\)
its function and parent set are unchanged, so its structural kernel is
the original \(k_i\). For \(i\in I\) its parent set is empty
and its function has constant value \(x_i\). Its new kernel is
\[
 \sum_{u\in\mathcal U_i}\mathbf1_{\{b=x_i\}}r_i(u)
       =\mathbf1_{\{b=x_i\}},
\]
since \(\sum_u r_i(u)=1\). Apply
\cref{thm:p16-causal-factorization} to this acyclic modified model.
The product contains one such indicator for each intervened coordinate
and the original factors at all other coordinates. If any intervened
coordinate disagrees with \(x\), its indicator makes the product zero.
If all agree, the indicators are one and the remaining product is exactly
the claimed expression, with those parent values fixed to \(x\).

The marginal formula is the grouping-of-fibers result already proved for
finite solution laws. When \(I=N\), every coordinate is fixed; the
remaining product is the empty product one, and the law has mass one
exactly at the single tuple \(x\). The proof uses the structural
kernels and never assigns a conditional ratio to a null event.
\end{proof}
\end{annexitem}

\begin{annexitem}{annex:p16.causal-positive-identification}{thm:p16-causal-positive-identification}{A positive observational law determines the kernels}
\begin{proof}
Let \(a\in S_{\operatorname{pa}_G(i)}\). All value sets are
nonempty and finite, so choose arbitrary values for the coordinates
outside the parent set to extend \(a\) to some \(w\in S_N\).
Full positive support gives \(p_M(w)>0\). Thus
\[
 \mathbb P(V_{\operatorname{pa}_G(i)}=a)
   =\sum_{z:z_{\operatorname{pa}_G(i)}=a}p_M(z)
   \geq p_M(w)>0.
\]
When the parent set is empty, the same probability is one. By
\cref{thm:p16-causal-factorization}, every kernel value is its
observational conditional ratio. Express its joint-event numerator and
parent-event denominator by their finite marginal sums. Substitution
into \cref{thm:p16-causal-truncation} gives the displayed product,
including the empty product when all coordinates are intervened upon.

Suppose now that two models in the specified class have equal full
observational laws. They have the same graph and value sets, so the
numerators and denominators in these ratios use the same index sets and
equal masses. They consequently have equal kernels and equal truncated
products at every assignment. Outside the intervention assignment both
laws are zero. This is equality of their full intervention laws and proves
identification by \cref{def:p16-causal-identification}.
Summing equal laws gives equal marginals; multiplying their masses by a
fixed finite response function and summing gives equal response means.
No inference about estimation from a finite sample enters this argument.
\end{proof}
\end{annexitem}

\begin{annexitem}{annex:p16.causal-confounded-example}{thm:p16-causal-confounded-example}{Computing the common-cause model}
\begin{proof}
The proposed outcome weights sum to
\[
 \sum_{z=0}^1\sum_{a=0}^1\sum_{b=0}^1
          \tfrac12\alpha(a)\alpha(b)
 =\Bigl(\sum_{z=0}^1\tfrac12\Bigr)
    \Bigl(\sum_{a=0}^1\alpha(a)\Bigr)
    \Bigl(\sum_{b=0}^1\alpha(b)\Bigr)=1.
\]
Their factorization gives the stated independent coordinate disturbances.
The graph's order \(Z,X,Y\) evaluates to
\(V(z,a,b)=(z,z\oplus a,z\oplus b)\).
For given \(z,x,y\), the unique disturbance coordinates giving
that assignment are \(a=x\oplus z\), \(b=y\oplus z\).
Consequently
\[
 p_{M_A}(z,x,y)=\tfrac12\alpha(x\oplus z)\alpha(y\oplus z).
\]
For \(x=y=0\), summing over \(z\) gives
\[
 \tfrac12\bigl((\tfrac34)^2+(\tfrac14)^2\bigr)=\tfrac5{16}.
\]
For \(x=y=1\) the two terms are reversed, giving the same value.
For either unequal pair, the sum is
\[
 \tfrac12\bigl((\tfrac34)(\tfrac14)
                    +(\tfrac14)(\tfrac34)\bigr)=\tfrac3{16}.
\]
This proves every table entry. Both rows and both columns sum to
\(8/16=1/2\). Dividing each \(Y=1\) entry by its positive row
sum gives \(5/8\) for \(X=1\) and \(3/8\) for \(X=0\).

Under \(\operatorname{do}(X=x)\), the equations instead evaluate
to \((Z,X,Y)=(z,x,z\oplus b)\). The outcome weights remain
\(\tfrac12\alpha(a)\alpha(b)\). The \(Y=1\) event consists
of \(z=0,b=1\) or \(z=1,b=0\), with arbitrary \(a\).
Its mass is
\[
 \tfrac12\alpha(1)\sum_a\alpha(a)
 +\tfrac12\alpha(0)\sum_a\alpha(a)
 =\tfrac18+\tfrac38=\tfrac12.
\]
This expression does not depend on \(x\). For a bit-valued response,
its expectation equals its probability of being one, by summing the
zero and one values. The observational difference is therefore
\(5/8-3/8=1/4\), and the intervention difference is
\(1/2-1/2=0\).

For the adjustment identity, sum the three-variable law over \(y\):
\[
 \mathbb P(X=x,Z=z)=\tfrac12\alpha(x\oplus z)>0.
\]
The numerator for \(Y=1\) is this same positive factor multiplied
by \(\alpha(1\oplus z)\). Dividing yields
\[
 \mathbb P(Y=1\mid X=x,Z=z)=\alpha(1\oplus z).
\]
It is \(1/4\) at \(z=0\) and \(3/4\) at \(z=1\).
Weighting by \(\mathbb P(Z=z)=1/2\) gives
\(1/8+3/8=1/2\), the intervention probability already calculated.
Every division here has an explicit positive denominator. This is a
calculation in the displayed model, with no appeal to a general graphical
adjustment theorem.
\end{proof}
\end{annexitem}

\begin{annexitem}{annex:p16.causal-nonidentification}{thm:p16-causal-nonidentification}{Two pairs of observationally equal models}
\begin{proof}
For the first assertion, define \(M_A'\) with endogenous variables
only \(X,Y\), each binary, and allowed graph \(X\to Y\).
Use the same three-bit outcome space and weights as \(M_A\).
Give its two disturbances value sets \(\{0,1\}^2\), with
\[
 \widetilde U_X(z,a,b)=(z,a),\qquad
 \widetilde U_Y(z,a,b)=(z,b).
\]
Define the structural functions by
\[
 f_X((),(z,a))=z\oplus a,\qquad
 f_Y(x,(z,b))=z\oplus b.
\]
Here \(()\) denotes the empty parent tuple. The second function is
total on its parent input \(x\) and disturbance, and ignores \(x\).
This is permitted by the model definition. Its graph is acyclic; its
disturbances need not be independent in the stated class. Evaluation
gives exactly the \(X,Y\) projection of \(M_A\), both before and
after replacing the equation for \(X\). Thus its observational law
is the displayed table and its target is \(1/2\).

For comparison, define \(M_B\) on outcome space \(\{0,1\}^2\)
with coordinates \((c,d)\) and weights
\[
 q_B(c,d)=\tfrac12\beta(d),\qquad
 \beta(0)=\tfrac58,\quad\beta(1)=\tfrac38.
\]
These weights are nonnegative and sum to
\((1/2+1/2)(5/8+3/8)=1\). Give it coordinate disturbances
\(U_X=c\), \(U_Y=d\), the same binary endogenous value sets,
the same graph \(X\to Y\), and equations
\[
 X=U_X,\qquad Y=X\oplus U_Y.
\]
The disturbances are independent; hence this model also belongs to the
larger class allowing dependence. Its solution is \((c,c\oplus d)\).
For a fixed observed pair \((x,y)\), the sole compatible outcome
is \((c,d)=(x,x\oplus y)\). Therefore
\[
 p_{M_B}(x,y)=\tfrac12\beta(x\oplus y)
  =\begin{cases}5/16,&x=y,\\3/16,&x\ne y.\end{cases}
\]
This agrees at all four pairs with \(p_{M_A'}\), and all four
masses are positive. Under \(\operatorname{do}(X=1)\), however,
\(Y=1\oplus d\), which is one exactly when \(d=0\). Its
probability is \(5/8\), unequal to \(1/2\). The two models
have the same specified graph, observed variables and observed law, but
different targets. They are the required nonidentification witness.

For the second assertion, use the same graph and binary endogenous value
sets, but let the outcome space be a singleton of mass one. Both
disturbance value sets are singletons and both disturbances are constant.
They are mutually independent: the sole joint tuple and each sole
marginal value have probability one. In \(M_C\) put \(X=0\)
and \(Y=0\), with the function for \(Y\) ignoring its allowed
parent. In \(M_D\) put \(X=0\) and \(Y=X\).
All functions are total on their full finite parent domains.
Both models evaluate observationally to \((0,0)\), so their laws
are equal point masses. Under \(\operatorname{do}(X=1)\), the
first gives \(Y=0\) and the second gives \(Y=1\). Their targets
are zero and one. The observational conditional on \(X=1\) is
undefined in both models, and the observation law supplies no value for
the response kernel at that parent input. This proves the claimed
nonidentification without the positive-support restriction.
\end{proof}
\end{annexitem}

\begin{annexitem}{annex:p15-w02c-p15-009-five-variable}{thm:p15-w02c-p15-009-five-variable}{Truncation, marginalization and positive conditioning events}
\begin{proof}
\paragraph{Item 1.} The parent sets read from the listed edges are
\(\operatorname{pa}(Z_1)=\operatorname{pa}(Z_2)=\varnothing\),
\(\operatorname{pa}(Z_3)=\{Z_1,Z_2\}\),
\(\operatorname{pa}(X)=\{Z_1,Z_3\}\) and
\(\operatorname{pa}(Y)=\{Z_2,Z_3,X\}\). The graph is acyclic, since
\(Z_1,Z_2,Z_3,X,Y\) is a topological order. The disturbances are
mutually independent, so \cref{thm:p16-causal-factorization} gives the
displayed product at every assignment.

\paragraph{Item 2.} Apply \cref{thm:p16-causal-truncation} with
\(I=\{X\}\) and value \(x_0\). The intervention law of
\((z_1,z_2,z_3,x,y)\) is zero unless \(x=x_0\), and for \(x=x_0\) it is
the product of the four kernels other than \(k_X\), with \(x_0\) in the
parent position of \(k_Y\). The marginal on \((Z_1,Z_2,Z_3,Y)\) sums
over \(x\); only the term \(x=x_0\) is nonzero, which gives the first
display. Summing it over \((z_1,z_2,z_3)\) gives the second.

\paragraph{Item 3.} Apply \cref{thm:p16-causal-truncation} with
\(I=\{X,Z_2\}\) and values \((x,z_2)\). The law is zero unless the
intervened coordinates equal \(x\) and \(z_2\), and otherwise it is the
product of \(k_1\), \(k_3\) and \(k_Y\), with \(z_2\) and \(x\) in their
parent positions. Summing over the two intervened coordinates leaves
the single nonzero term, and summing over \((z_1,z_3)\) gives the
second display.

\paragraph{Item 4: marginals of the observed law.} Sum the product of
item~1 over \(y\) and then over \(x\). Since
\(\sum_yk_Y(y\mid z_2,z_3,x)=1\) and \(\sum_xk_X(x\mid z_1,z_3)=1\),
\[
 \mathbb P(Z_1=z_1,Z_2=z_2,Z_3=z_3)
   =k_1(z_1)k_2(z_2)k_3(z_3\mid z_1,z_2),
\]
and a further sum over \(z_3\) gives
\(\mathbb P(Z_1=z_1,Z_2=z_2)=k_1(z_1)k_2(z_2)\). For a vertex with an
empty parent set the parent event has probability one, so
\cref{thm:p16-causal-factorization} gives \(k_1(z_1)=\mathbb P(Z_1=z_1)\)
and \(k_2(z_2)=\mathbb P(Z_2=z_2)\) at every value.

\paragraph{Item 4: the single intervention.} In the last display of
item~2, a term with \((z_1,z_2,z_3)\notin\mathcal Z^+\) has
\(k_1(z_1)k_2(z_2)k_3(z_3\mid z_1,z_2)=0\) by the marginal just
computed, so it vanishes. Fix \((z_1,z_2,z_3)\in\mathcal Z^+\). Then
\(k_1(z_1)k_2(z_2)>0\), so \(\mathbb P(Z_1=z_1,Z_2=z_2)>0\), and
\cref{thm:p16-causal-factorization} gives
\(k_3(z_3\mid z_1,z_2)=\mathbb P(z_3\mid z_1,z_2)\). By hypothesis
\(\mathbb P(Z_2=z_2,Z_3=z_3,X=x_0)>0\), and the same theorem gives
\(k_Y(y\mid z_2,z_3,x_0)=\mathbb P(y\mid z_2,z_3,x_0)\). Substituting
these and the root identities into the remaining terms gives the first
observational formula.

\paragraph{Item 4: the simultaneous intervention.} In the last display
of item~3, a term with \(k_1(z_1)k_3(z_3\mid z_1,z_2)=0\) vanishes. If
\(k_1(z_1)k_3(z_3\mid z_1,z_2)>0\), then, because
\(k_2(z_2)=\mathbb P(Z_2=z_2)>0\) by hypothesis,
\(\mathbb P(Z_1=z_1,Z_2=z_2,Z_3=z_3)>0\), that is
\((z_1,z_2,z_3)\in\mathcal Z^+\). Conversely a triple in
\(\mathcal Z^+\) has \(k_1(z_1)k_3(z_3\mid z_1,z_2)>0\). So the sum may
be restricted to the pairs \((z_1,z_3)\) with
\((z_1,z_2,z_3)\in\mathcal Z^+\). For such a pair,
\(\mathbb P(Z_1=z_1,Z_2=z_2)>0\) gives
\(k_3(z_3\mid z_1,z_2)=\mathbb P(z_3\mid z_1,z_2)\), the hypothesis
\(\mathbb P(Z_2=z_2,Z_3=z_3,X=x)>0\) gives
\(k_Y(y\mid z_2,z_3,x)=\mathbb P(y\mid z_2,z_3,x)\), and
\(k_1(z_1)=\mathbb P(Z_1=z_1)\). Substitution gives the second
observational formula. Every conditional probability used was
introduced only at a conditioning event already shown to have positive
probability.
\end{proof}
\end{annexitem}

\begin{annexitem}{annex:p15-w02c-p15-009-chain-weak}{thm:p15-w02c-p15-009-chain-weak}{Independence of one disturbance from the treatment's background}
\begin{proof}
By topological evaluation (\cref{thm:p16-causal-solution}),
\(Z=f_Z(U_Z)\), \(X=f_X(Z,U_X)=g(U_Z,U_X)\) with
\(g(u_Z,u_X)=f_X(f_Z(u_Z),u_X)\), and \(Y=f_Y(X,U_Y)\) at every
outcome. In the model \(M_{\{X\}=x}\) the equations of \(Z\) and \(Y\)
are unchanged and \(X\) is the constant \(x\), so its solution is
\(Z^{X=x}=f_Z(U_Z)=Z\), \(X^{X=x}=x\) and \(Y^{X=x}=f_Y(x,U_Y)\).

Fix \(y\) and put \(A=\{u:f_Y(x,u)=y\}\) and
\(W=\{w:g(w)=x\}\), sets of values of \(U_Y\) and of \((U_Z,U_X)\). Then
\(\mathbb P_M(Y=y\mid\operatorname{do}(X=x))=\mathbb P(U_Y\in A)\). On
the event \(\{X=x\}\) the equation for \(Y\) reads \(Y=f_Y(x,U_Y)\),
so \(\{Y=y,X=x\}=\{U_Y\in A\}\cap\{(U_Z,U_X)\in W\}\). Summing the
hypothesized product identity over \(u\in A\) and \(w\in W\) gives
\[
 \mathbb P(Y=y,X=x)=\mathbb P(U_Y\in A)\,\mathbb P((U_Z,U_X)\in W)
                   =\mathbb P(U_Y\in A)\,\mathbb P(X=x).
\]
Dividing by \(\mathbb P(X=x)>0\) proves the first identity, with common
value \(\mathbb P(f_Y(x,U_Y)=y)\). Multiplying by \(h(y)\) and summing
over the finite set \(S_Y\) gives
\(\mathbb E[h(Y^{X=x})]=\sum_yh(y)\mathbb P(Y=y\mid X=x)
=\mathbb E[h(Y)\mid X=x]\).

For the joint law, put \(W_z=\{w=(u_Z,u_X):f_Z(u_Z)=z\}\). Then
\(\{Z^{X=x}=z,Y^{X=x}=y\}=\{(U_Z,U_X)\in W_z\}\cap\{U_Y\in A\}\), and
the same summation gives
\(\mathbb P((U_Z,U_X)\in W_z)\,\mathbb P(U_Y\in A)
=\mathbb P(Z=z)\,\mathbb P(Y=y\mid X=x)\). The independence of \(U_Z\)
from \(U_X\) was never used, and \(f_Y\) entered only through its
arguments \(X\) and \(U_Y\); these are exactly the two assumptions the
source names.
\end{proof}
\end{annexitem}

\begin{annexitem}{annex:p15-w02c-p15-009-linear-effect}{thm:p15-w02c-p15-009-linear-effect}{Means and covariances under a linear response equation}
\begin{proof}
Item~1. In \(M_{\{X\}=x_0}\) the equation of \(Y\), its parent set and
its disturbance are unchanged, and \(X^{X=x_0}=x_0\) at every outcome
(\cref{def:p16-causal-intervention}). Since the only parent of \(Y\) is
\(X\), the solution satisfies
\(Y^{X=x_0}(\omega)=f_Y(x_0,U_Y(\omega))=\beta x_0+U_Y(\omega)\) for
every \(\omega\). The outcome weights \(q\) are unchanged, so linearity
of the finite expectation gives
\(\mathbb E[Y^{X=x_0}]=\beta x_0+\mathbb E[U_Y]\). Subtracting the value
at \(x_0\) from that at \(x_0'\) gives \(\beta(x_0'-x_0)\). No property
of \(f_X\), of the other equations or of the joint law of the
disturbances was used.

Item~2. In the unmodified model \(Y(\omega)=\beta X(\omega)+U_Y(\omega)\)
at every outcome. If \(\mathbb P(X=x)>0\), then
\[
 \mathbb E[Y\mid X=x]
 =\frac{\sum_{\omega:X(\omega)=x}q(\omega)\bigl(\beta x+U_Y(\omega)\bigr)}
       {\mathbb P(X=x)}
 =\beta x+\mathbb E[U_Y\mid X=x].
\]
Comparing with item~1 at \(x_0=x\), the two means are equal exactly when
\(\mathbb E[U_Y\mid X=x]=\mathbb E[U_Y]\). If
\(\mathbb P(U_Y=u,X=x)=\mathbb P(U_Y=u)\mathbb P(X=x)\) for all \(u\),
then
\(\mathbb E[U_Y\mid X=x]=\sum_uu\,\mathbb P(U_Y=u,X=x)/\mathbb P(X=x)
=\sum_uu\,\mathbb P(U_Y=u)=\mathbb E[U_Y]\).

Item~3. Covariance is bilinear and \(Y=\beta X+U_Y\) pointwise, so
\(\operatorname{Cov}(X,Y)=\beta\operatorname{Cov}(X,X)
+\operatorname{Cov}(X,U_Y)=\beta\operatorname{Var}(X)
+\operatorname{Cov}(X,U_Y)\). Dividing by \(\operatorname{Var}(X)>0\),
the slope equals \(\beta\) exactly when \(\operatorname{Cov}(X,U_Y)=0\).
If \(X=U_X\), substitute \(U_X\) for \(X\) in the covariance term.

Item~4. In the first model take \(\mathcal U_X=\{-1,1\}\),
\(\mathcal U_Y=\{-\gamma,\gamma\}\), \(f_X(u)=u\), the graph \(X\to Y\),
and \(S_Y=\{\beta x+u:x\in\mathcal U_X,u\in\mathcal U_Y\}\); the
functions are total and the disturbances are dependent, which the model
definition permits. Here \(\mathbb E[U_Y]=\gamma(\tfrac12-\tfrac12)=0\),
so item~1 gives \(\mathbb E[Y^{X=x}]=\beta x\). The event \(\{X=x\}\) is
the single outcome \(\omega=x\), of mass \(\tfrac12\), on which
\(U_Y=\gamma x\); item~2 gives
\(\mathbb E[Y\mid X=x]=\beta x+\gamma x\), which differs from \(\beta x\)
since \(\gamma\ne0\) and \(x\ne0\).
In the second model take \(\mathcal U_X=\{-1,0,1\}\),
\(\mathcal U_Y=\{\tfrac13,-\tfrac23\}\), \(f_X(u)=u\), the same graph and
\(S_Y\) as before. Then \(\mathbb E[X]=0\),
\(\mathbb E[U_Y]=\tfrac13(1+0+1)-\tfrac23=0\), and
\(\operatorname{Cov}(X,U_Y)=\mathbb E[XU_Y]
=\tfrac13\bigl((-1)\tfrac13+0+1\cdot\tfrac13\bigr)=0\), while
\(\operatorname{Var}(X)=\tfrac23>0\); by item~3 the slope is \(\beta\).
Item~1 gives \(\mathbb E[Y^{X=0}]=0\). The event \(\{X=0\}\) is the
single outcome \(\omega=0\), on which \(U_Y=-\tfrac23\), so item~2 gives
\(\mathbb E[Y\mid X=0]=-\tfrac23\).

Finally, under the standardization \(\mathbb E[X]=\mathbb E[Y]=0\),
taking expectations in \(Y=\beta X+U_Y\) gives \(\mathbb E[U_Y]=0\), and
item~1 reads \(\mathbb E[Y^{X=x_0}]=\beta x_0\). With
\(\operatorname{Var}(X)=1\) and \(X=U_X\), item~3 reads
\(\operatorname{Cov}(X,Y)=\beta+\operatorname{Cov}(U_X,U_Y)\), which is
\(\beta\) when the disturbance covariance vanishes.
\end{proof}
\end{annexitem}

\begin{annexitem}{annex:p15-w02c-p15-009-diagram-paths}{thm:p15-w02c-p15-009-diagram-paths}{Kinship and paths in the diagram}
\begin{proof}
Item~1. By \cref{def:p15-w02c-p15-009-causal-diagram}, a solid arrow
joins two endogenous vertices, a dashed arrow \(U_i\to i\) ends at an
endogenous vertex, and a double arrow is not an arrow. Hence no arrow
ends at a disturbance vertex, so \(U_i\) has no parent. The only arrow
that starts at \(U_i\) is \(U_i\to i\), so \(i\) is its only child.

Let \(U_i=a_0\to a_1\to\cdots\to a_m\), with \(m\geq1\), be a directed
path. Its first arrow is \(U_i\to i\), so \(a_1=i\). Every later arrow
starts at an endogenous vertex, because no arrow ends at a disturbance
vertex, and an arrow starting at an endogenous vertex is solid; so it
belongs to \(E\) and ends at an endogenous vertex. Thus
\(a_1,\ldots,a_m\) is a directed path in \(G\) from \(i\), of length
\(m-1\geq0\), and \(a_m\in\{i\}\cup\operatorname{de}_G(\{i\})\).
Conversely, \(U_i\to i\) followed by a directed path of \(G\) from \(i\)
is a directed path of the diagram. This identifies the descendants of
\(U_i\).

A directed cycle cannot pass through a disturbance vertex, since it would
have to enter that vertex by an arrow. It would therefore use only
endogenous vertices and, by the previous paragraph's observation, only
solid arrows, so it would be a directed cycle of \(G\), which
\cref{def:p16-causal-dag} excludes. The same observation shows that a
directed path between two endogenous vertices uses only solid arrows, so
it is a directed path of \(G\); and an arrow between two endogenous
vertices is an edge of \(E\). Hence, for endogenous \(i,j\), the
statements that \(j\) is a parent, child, ancestor or descendant of \(i\)
mean the same in the diagram as in \(G\). The disturbance vertices that
are ancestors of an endogenous \(i\) are, by the identification of their
descendants, the \(U_k\) with \(i\in\{k\}\cup\operatorname{de}_G(\{k\})\).

Item~2. In the diagram of Fig.~1(b) the edges are \(X\to Y\),
\(U_X\to X\), \(U_Y\to Y\) and \(U_X\leftrightarrow U_Y\). Let
\(w_0=X,\ldots,w_m=Y\) be a path. Since its vertices are distinct and
\(w_m=Y\), the vertex \(Y\) occurs only at the end. The edges at \(X\)
join it to \(Y\) and to \(U_X\). If \(w_1=Y\), then \(m=1\) and the path
is \((X,Y)\). Otherwise \(w_1=U_X\). The edges at \(U_X\) join it to
\(X\) and to \(U_Y\); \(X\) is already used, so \(w_2=U_Y\). The edges at
\(U_Y\) join it to \(U_X\) and to \(Y\); \(U_X\) is already used, so
\(w_3=Y\), \(m=3\), and the path is \((X,U_X,U_Y,Y)\). Both sequences
consist of distinct vertices with consecutive ones adjacent, so both are
paths, and there are no others. In the diagram of Fig.~1(a) the double
arrow is absent, the only edge at \(U_X\) joins it to \(X\), and the
second case stops at \(w_1=U_X\) without reaching \(Y\); the only path is
\((X,Y)\).
\end{proof}
\end{annexitem}

\begin{annexitem}{annex:p15-w02c-p15-009-path-solve}{thm:p15-w02c-p15-009-path-solve}{The covariance solution and three observationally equal models}
\begin{proof}
Item~1. A Fig.~1(a) model satisfies the hypotheses of
\cref{thm:p15-w02c-p15-009-linear-effect}, with \(X\) parentless and
\(X=U_X\). Item~3 of that theorem gives
\(\operatorname{Cov}(X,Y)=\beta\operatorname{Var}(X)
+\operatorname{Cov}(U_X,U_Y)=\beta\operatorname{Var}(X)\), and dividing
by \(\operatorname{Var}(X)>0\) gives the displayed identity. Item~1 of
that theorem gives
\(\mathbb E[Y^{X=x'}]-\mathbb E[Y^{X=x}]=\beta(x'-x)\), and substituting
the identity gives the stated form. With \(\mu(x,y)=\mathbb P(X=x,Y=y)\),
\[
\begin{aligned}
 \operatorname{Var}(X)&=\sum_{x,y}x^2\mu(x,y)
     -\Bigl(\sum_{x,y}x\,\mu(x,y)\Bigr)^2,\\
 \operatorname{Cov}(X,Y)&=\sum_{x,y}xy\,\mu(x,y)
     -\sum_{x,y}x\,\mu(x,y)\sum_{x,y}y\,\mu(x,y),
\end{aligned}
\]
so both depend on the model only through \(\mu\). Two Fig.~1(a) models
with the same \(\mu\) and \(\operatorname{Var}(X)>0\) therefore have the
same \(\beta\), and the same effect at any \(x,x'\) lying in both value
sets. If \(\operatorname{Var}(X)=1\), the identity reads
\(\beta=\operatorname{Cov}(X,Y)\).

Item~2. Take \(S_X=S_Y=\{-1,0,1\}\) in all three models, and
disturbance value sets \(\mathcal U_X=\{-1,1\},\mathcal U_Y=\{0\}\) in
\(M_1\), \(\mathcal U_X=\mathcal U_Y=\{-1,1\}\) in \(M_2\), and
\(\mathcal U_X=\{0\},\mathcal U_Y=\{-1,1\}\) in \(M_3\). Every structural
function is total with values in \(\{-1,0,1\}\). Evaluating in
topological order: in \(M_1\), \(X=\omega\) and \(Y=\omega+0=\omega\);
in \(M_2\), \(X=\omega\) and \(Y=0\cdot\omega+\omega=\omega\); in
\(M_3\), \(Y=\omega\) and \(X=\omega+0=\omega\). Each model therefore
puts mass \(\tfrac12\) on \((-1,-1)\) and on \((1,1)\). For this law
\(\mathbb E[X]=\mathbb E[Y]=0\) and
\(\operatorname{Var}(X)=\operatorname{Cov}(X,Y)=\mathbb E[XY]=1\).

In \(M_1\) the disturbance \(U_Y\) is constant, so
\(\operatorname{Cov}(U_X,U_Y)=0\), and \(M_1\) is a Fig.~1(a) model with
coefficient \(1\); by \cref{def:p16-causal-intervention} its intervened
response is \(Y^{X=x}=x+U_Y=x\), whose means at \(x=1\) and \(x=-1\)
differ by \(2\). In \(M_2\),
\(\operatorname{Cov}(U_X,U_Y)=\mathbb E[\omega^2]-0=1\); the intervened
response is \(Y^{X=x}=0\cdot x+U_Y=\omega\), with mean \(0\) for each
\(x\). In \(M_3\), \(U_X\) is constant, so
\(\mathbb P(U_X=0,U_Y=u)=\mathbb P(U_Y=u)
=\mathbb P(U_X=0)\,\mathbb P(U_Y=u)\) and the disturbances are mutually
independent. The intervention \(\operatorname{do}(X=x)\) replaces the
equation of \(X\) and leaves the parentless equation \(Y=U_Y\)
unchanged, so \(Y^{X=x}=\omega\), with mean \(0\) for each \(x\). Hence
the intervention contrast is \(2\) in \(M_1\) and \(0\) in \(M_2\) and
\(M_3\), while the common law gives
\(\operatorname{Cov}(X,Y)/\operatorname{Var}(X)=1\) in all three.
\end{proof}
\end{annexitem}

\begin{annexitem}{annex:p15-w02c-p15-009-chain-testable}{thm:p15-w02c-p15-009-chain-testable}{Testing the chain}
\begin{proof}
Item~1. By \cref{thm:p16-causal-factorization}, with the parent sets
\(\varnothing,\{Z\},\{X\}\),
\[
 \mathbb P(Z=z,X=x,Y=y)=k_Z(z)\,k_X(x\mid z)\,k_Y(y\mid x).
\]
Summing over \(y\) and using \(\sum_yk_Y(y\mid x)=1\) gives
\(\mathbb P(Z=z,X=x)=k_Z(z)k_X(x\mid z)\). Summing the product over
\(z\) gives \(\mathbb P(X=x,Y=y)=m(x)k_Y(y\mid x)\), where
\(m(x)=\sum_zk_Z(z)k_X(x\mid z)\); summing this over \(y\) shows
\(m(x)=\mathbb P(X=x)\). Hence
\[
\begin{aligned}
 \mathbb P(Z=z,X=x,Y=y)\,\mathbb P(X=x)
 &=k_Z(z)k_X(x\mid z)\cdot m(x)k_Y(y\mid x)\\
 &=\mathbb P(Z=z,X=x)\,\mathbb P(X=x,Y=y).
\end{aligned}
\]
Dividing by \(\mathbb P(X=x)^2>0\) gives the identity of item~1.

Item~2. Choose \(z_0\ne z_1\) in \(S_Z\), \(y_0\ne y_1\) in \(S_Y\)
and \(x_0\in S_X\), and let \(p\) put mass \(\tfrac12\) on
\((z_0,x_0,y_0)\) and on \((z_1,x_0,y_1)\). Then
\(\mathbb P(X=x_0)=1\),
\(\mathbb P(Z=z_0,Y=y_0\mid X=x_0)=\tfrac12\), and
\(\mathbb P(Z=z_0\mid X=x_0)\,\mathbb P(Y=y_0\mid X=x_0)
=\tfrac12\cdot\tfrac12=\tfrac14\). By item~1 no model of that item has
\(p\) as its observational law.

Item~3, the model \(M_{\mathrm a}\). Let
\(\Omega=S_Z\times S_X\times S_Y\) with \(q=p\), and let each
disturbance be the identity map of \(\Omega\), so
\(\mathcal U_Z=\mathcal U_X=\mathcal U_Y=\Omega\). For
\(\omega=(\omega_Z,\omega_X,\omega_Y)\) put
\(f_Z(\omega)=\omega_Z\), \(f_X(z,\omega)=\omega_X\) and
\(f_Y(x,\omega)=\omega_Y\). These are total functions with the parent
sets of the graph \(Z\to X\to Y\), and evaluation in that order gives
\((Z,X,Y)(\omega)=\omega\). The observational law is therefore
\(\mathbb P((Z,X,Y)=\omega)=q(\omega)=p(\omega)\).

Item~3, the model \(M_{\mathrm b}\). Put
\(p_Z(z)=\sum_{x,y}p(z,x,y)\) and \(p_{ZX}(z,x)=\sum_yp(z,x,y)\), fix
\(x^*\in S_X\) and \(y^*\in S_Y\), and define
\[
 \kappa_X(x\mid z)=
 \begin{cases}p_{ZX}(z,x)/p_Z(z),&p_Z(z)>0,\\
  \mathbf 1[x=x^*],&p_Z(z)=0,\end{cases}
\]
\[
 \kappa_Y(y\mid z,x)=
 \begin{cases}p(z,x,y)/p_{ZX}(z,x),&p_{ZX}(z,x)>0,\\
  \mathbf 1[y=y^*],&p_{ZX}(z,x)=0.\end{cases}
\]
For each fixed second argument, each is nonnegative and sums to one over
its first argument. Let \(\mathcal U_Z=S_Z\), let \(\mathcal U_X\) be the
finite set of functions \(\varphi:S_Z\to S_X\), and let \(\mathcal U_Y\)
be the finite set of functions \(\psi:S_Z\times S_X\to S_Y\). Put
\(\Omega=\mathcal U_Z\times\mathcal U_X\times\mathcal U_Y\) with
\(q(u,\varphi,\psi)=r_Z(u)\,r_X(\varphi)\,r_Y(\psi)\), where
\[
\begin{aligned}
 r_Z(u)&=p_Z(u),\qquad
 r_X(\varphi)=\prod_{z\in S_Z}\kappa_X(\varphi(z)\mid z),\\
 r_Y(\psi)&=\prod_{(z,x)\in S_Z\times S_X}\kappa_Y(\psi(z,x)\mid z,x),
\end{aligned}
\]
and let \(U_Z,U_X,U_Y\) be the three coordinate maps. By finite
distributivity, \(\sum_\varphi r_X(\varphi)
=\prod_z\sum_x\kappa_X(x\mid z)=1\), and likewise
\(\sum_\psi r_Y(\psi)=1\) and \(\sum_ur_Z(u)=1\). Hence \(q\) is a mass
function, the three factors are the laws of the coordinates, and the
disturbances are mutually independent. Take the graph with edges
\(Z\to X\), \(X\to Y\), \(Z\to Y\) and the structural functions
\(f_Z(u)=u\), \(f_X(z,\varphi)=\varphi(z)\) and
\(f_Y((z,x),\psi)=\psi(z,x)\). The structural kernels of
\cref{def:p16-causal-kernels} are \(k_Z=p_Z\) and
\[
 k_X(x\mid z)=\sum_{\varphi:\varphi(z)=x}r_X(\varphi)
 =\kappa_X(x\mid z)\prod_{z'\ne z}\sum_{x'\in S_X}\kappa_X(x'\mid z')
 =\kappa_X(x\mid z),
\]
and in the same way \(k_Y(y\mid z,x)=\kappa_Y(y\mid z,x)\). By
\cref{thm:p16-causal-factorization} the observational law is
\(p_Z(z)\kappa_X(x\mid z)\kappa_Y(y\mid z,x)\). If
\(p_{ZX}(z,x)>0\), then \(p_Z(z)>0\) and this product is
\(p_Z(z)\cdot\frac{p_{ZX}(z,x)}{p_Z(z)}\cdot
\frac{p(z,x,y)}{p_{ZX}(z,x)}=p(z,x,y)\). If \(p_{ZX}(z,x)=0\), then
\(p(z,x,y)=0\); either \(p_Z(z)=0\), or \(p_Z(z)>0\) and
\(\kappa_X(x\mid z)=0\), and in both cases the product is \(0\). So the
observational law is \(p\).

The four missing arrows \(X\to Z\), \(Y\to X\), \(Z\to Y\), \(Y\to Z\)
are absent from the graph of \(M_{\mathrm a}\), and the three missing
double arrows hold in \(M_{\mathrm b}\), whose disturbances are mutually
independent. Since both models have the arbitrary law \(p\), no law is
excluded by any one assumption alone.

Item~4. Under equal masses on \(\{0,1\}^2\), each of \(U_Z=a\),
\(U_X=b\) and \(U_Y=a\oplus b\) is a fair bit. For each pair of them the
outcome is determined by the pair's values (from \(a\) and \(b\)
directly, and from either of them together with \(a\oplus b\) by one
further sum modulo two), so the four value pairs correspond one to one to
the four outcomes; each pair has mass \(\tfrac14\) at every value pair
and is therefore independent. Topological evaluation gives \(Z=a\),
\(X=b\) and \(Y=b\oplus(a\oplus b)=a\). The event \(\{X=0\}\) consists
of \((0,0)\) and \((1,0)\), of total mass \(\tfrac12\); on it
\(Z=Y=a\). Hence \(\mathbb P(Z=1,Y=1\mid X=0)=\tfrac12\), while
\(\mathbb P(Z=1\mid X=0)=\mathbb P(Y=1\mid X=0)=\tfrac12\), whose
product is \(\tfrac14\).
\end{proof}
\end{annexitem}

\begin{annexitem}{annex:p15-w02c-p15-009-chain-diagram}{thm:p15-w02c-p15-009-chain-diagram}{Separation in the chain's diagram and the converse of its test}
\begin{proof}
Item~1. By \cref{thm:p16-causal-solution}, at every outcome \(\omega\)
the solution satisfies \(Z(\omega)=f_Z(U_Z(\omega))\),
\(X(\omega)=f_X(Z(\omega),U_X(\omega))\) and
\(Y(\omega)=f_Y(X(\omega),U_Y(\omega))\); the parent set of \(Z\) in
\(G\) is empty, so \(f_Z\) has the disturbance as its only argument.
Fix \(t=(u_Z,u_X,u_Y,z,x,y)\). If \(z=f_Z(u_Z)\),
\(x=f_X(z,u_X)\) and \(y=f_Y(x,u_Y)\), the event that the six-tuple
equals \(t\) is the event \(\{(U_Z,U_X,U_Y)=(u_Z,u_X,u_Y)\}\), because
on that event the three equations force \(Z=z\), \(X=x\), \(Y=y\). If one
of the three equalities fails, the event is empty. By mutual independence
(\cref{def:p16-causal-kernels}) the probability of the first event is
\(r_Z(u_Z)r_X(u_X)r_Y(u_Y)\), which gives the displayed mass function.

In \(D\) the parent sets are \(\varnothing\) for each disturbance vertex,
\(\{U_Z\}\) for \(Z\), \(\{Z,U_X\}\) for \(X\) and \(\{X,U_Y\}\) for
\(Y\), by \cref{def:p15-w02c-p15-009-causal-diagram}. The factor
\(r_i\) of a disturbance vertex is a mass function on
\(\mathcal U_i\). For each value of its parents, the factor
\(\mathbf 1[z=f_Z(u_Z)]\) sums to one over \(z\in S_Z\), since exactly
one \(z\) equals \(f_Z(u_Z)\); the same holds for the factors of \(X\)
and \(Y\). The mass function therefore has the product form of
\cref{thm:p16-ancestral-marginal} for the DAG \(D\), which is acyclic by
item~1 of \cref{thm:p15-w02c-p15-009-diagram-paths}, with finite
nonempty value sets. \Cref{thm:p16-finite-global-markov} then gives, for
pairwise disjoint vertex sets \(A,B,C\) of \(D\) such that \(C\)
d-separates \(A\) from \(B\), the conditional independence of the
corresponding coordinate tuples under this mass function. Every
probability of an event determined by the six variables equals the mass
of the corresponding set of tuples, so the same identities hold for the
random variables themselves.

Item~2. The edges of \(D\), ignoring direction, join \(U_Z\) with \(Z\),
\(Z\) with \(X\), \(X\) with \(Y\), \(U_X\) with \(X\), and \(U_Y\)
with \(Y\). A path from \(U_Z\) must begin \(U_Z,Z\), since \(Z\) is the
only neighbour of \(U_Z\), and then continue to \(X\), the only
neighbour of \(Z\) besides \(U_Z\). From \(X\) the unused neighbours are
\(U_X\) and \(Y\). From \(U_X\) no unused neighbour remains, and from
\(Y\) the only unused neighbour is \(U_Y\), which has no other
neighbour. Hence the only path from \(U_Z\) to \(Y\) is
\((U_Z,Z,X,Y)\), and the only path from \(U_Z\) to \(U_X\) is
\((U_Z,Z,X,U_X)\).

On \((U_Z,Z,X,Y)\) the edges are \(U_Z\to Z\), \(Z\to X\), \(X\to Y\),
so the internal vertices \(Z\) and \(X\) are noncolliders. Given
\(\{Z\}\) the noncollider \(Z\) lies in the conditioning set, and given
\(\{X\}\) the noncollider \(X\) does; in both cases the only path is
blocked (\cref{def:p16-active-path}), which proves the two
d-separations. On \((U_Z,Z,X,U_X)\) the edges are \(U_Z\to Z\),
\(Z\to X\), \(X\leftarrow U_X\), so \(Z\) is a noncollider and \(X\) is
a collider. Given \(\varnothing\), the collider \(X\) is not in
\(\operatorname{an}_D(\varnothing)=\varnothing\), so the path is blocked
and \(\{U_Z\}\), \(\{U_X\}\) are d-separated. Given \(\{Y\}\), the
noncollider \(Z\) is outside \(\{Y\}\), and
\(X\in\operatorname{an}_D(\{Y\})\) because of the arrow \(X\to Y\); the
path is active and the two vertices are not d-separated. Finally,
\((U_X,X,Y)\) is a path with edges \(U_X\to X\to Y\); its internal
vertex \(X\) is a noncollider outside \(\{Z\}\), so the path is active
given \(\{Z\}\), and \(\{U_X\}\), \(\{Y\}\) are not d-separated given
\(\{Z\}\). The three displayed conditional independencies follow from
the three d-separations by item~1; the case of the empty conditioning set
is ordinary independence, as stated in
\cref{def:p16-graphical-conditional-independence}.

Item~3. The disturbance law is
\(\mathbb P(U_Z=0,U_X=u,U_Y=0)=\mathbb P(U_X=u)\), the product of the
three marginal mass functions, since \(U_Z\) and \(U_Y\) take the single
value \(0\); so the disturbances are mutually independent. The
functions are total and use only the declared parents. Evaluation gives
\(Z=0\), \(X=\omega\) and \(Y=\omega\) at each outcome. Hence
\(\mathbb P(Z=0)=1\), conditioning on \(Z=0\) changes no probability,
and \(\mathbb P(U_X=1,Y=1\mid Z=0)=q(1)=\tfrac12\), while
\(\mathbb P(U_X=1\mid Z=0)=\mathbb P(Y=1\mid Z=0)=\tfrac12\), whose
product is \(\tfrac14\).

Item~4. The four outcomes \((a,b)\) have mass \(\tfrac14\), so
\(\mathbb P(U_Z=s,U_X=t,U_Y=0)=\tfrac14=\tfrac12\cdot\tfrac12\cdot1\)
for all bits \(s,t\): the disturbances are mutually independent, and in
particular \(U_Z\) and \(U_X\) are independent. Evaluation gives
\(Z=a\) and \(X=Y=\max(a,b)\). The event \(\{Y=1\}\) consists of
\((0,1)\), \((1,0)\) and \((1,1)\), of mass \(\tfrac34\). Within it,
\(\{U_Z=1,U_X=1\}\) is \(\{(1,1)\}\), of mass \(\tfrac14\), so the
conditional probability is \(\tfrac13\); and \(\{U_Z=1\}\) is
\(\{(1,0),(1,1)\}\), of mass \(\tfrac12\), so
\(\mathbb P(U_Z=1\mid Y=1)=\tfrac23\). Exchanging the roles of \(a\) and
\(b\) gives \(\mathbb P(U_X=1\mid Y=1)=\tfrac23\). The product is
\(\tfrac49\ne\tfrac13\).

With \(f_X(z,u)=u\) instead, evaluation gives \(Y=X=b\). For a bit
\(c\) the event \(\{Y=c\}\) is \(\{(0,c),(1,c)\}\), of mass \(\tfrac12\).
For bits \(s,t\), the event \(\{U_Z=s,U_X=t,Y=c\}\) is \(\{(s,c)\}\)
when \(t=c\) and is empty otherwise, so
\(\mathbb P(U_Z=s,U_X=t\mid Y=c)=\tfrac12\mathbf 1[t=c]\). Also
\(\mathbb P(U_Z=s\mid Y=c)=\tfrac12\) and
\(\mathbb P(U_X=t\mid Y=c)=\mathbf 1[t=c]\), whose product is the same
number. Hence \(U_Z\mathrel{\perp\!\!\!\perp}U_X\mid Y\).

Item~5, necessity. If \(p\) is the observational law of such a model,
item~1 of \cref{thm:p15-w02c-p15-009-chain-testable} gives the
conditional identity at every \(x\) with \(p_X(x)>0\). Each conditional
probability there is a two- or three-coordinate mass divided by
\(p_X(x)\), so multiplying the identity by \(p_X(x)^2\) gives the
displayed identity.

Item~5, sufficiency. Suppose the identity holds. Let \(p_Z\) be the
\(Z\)-marginal, \(p_Z(z)=\sum_{x,y}p(z,x,y)\); fix \(x^*\in S_X\) and
\(y^*\in S_Y\); and define
\[
 \kappa_X(x\mid z)=
 \begin{cases}p_{ZX}(z,x)/p_Z(z),&p_Z(z)>0,\\
  \mathbf 1[x=x^*],&p_Z(z)=0,\end{cases}
\]
\[
 \kappa_Y(y\mid x)=
 \begin{cases}p_{XY}(x,y)/p_X(x),&p_X(x)>0,\\
  \mathbf 1[y=y^*],&p_X(x)=0.\end{cases}
\]
For each fixed second argument, each is nonnegative and sums to one over
its first argument. Let \(\mathcal U_Z=S_Z\), let \(\mathcal U_X\) be the
finite set of functions \(\varphi:S_Z\to S_X\), and let \(\mathcal U_Y\)
be the finite set of functions \(\psi:S_X\to S_Y\). Put
\(\Omega=\mathcal U_Z\times\mathcal U_X\times\mathcal U_Y\) with
\(q(u,\varphi,\psi)=r_Z(u)r_X(\varphi)r_Y(\psi)\), where
\[
 r_Z(u)=p_Z(u),\qquad
 r_X(\varphi)=\prod_{z\in S_Z}\kappa_X(\varphi(z)\mid z),
\]
\[
 r_Y(\psi)=\prod_{x'\in S_X}\kappa_Y(\psi(x')\mid x'),
\]
and let \(U_Z,U_X,U_Y\) be the coordinate maps. By finite
distributivity,
\[
 \sum_\varphi r_X(\varphi)=\prod_z\sum_{x}\kappa_X(x\mid z)=1,
\]
and likewise \(\sum_\psi r_Y(\psi)=1\) and \(\sum_ur_Z(u)=1\). Hence
\(q\) is a mass function whose three factors are the laws of the
coordinates, and the disturbances are mutually independent. Take the
graph \(Z\to X\to Y\) and \(f_Z(u)=u\), \(f_X(z,\varphi)=\varphi(z)\),
\(f_Y(x,\psi)=\psi(x)\). The structural kernels of
\cref{def:p16-causal-kernels} are \(k_Z=p_Z\),
\[
 k_X(x\mid z)=\sum_{\varphi:\varphi(z)=x}r_X(\varphi)
 =\kappa_X(x\mid z)\prod_{z'\ne z}\sum_{x'\in S_X}\kappa_X(x'\mid z')
 =\kappa_X(x\mid z),
\]
and, in the same way,
\(k_Y(y\mid x)=\kappa_Y(y\mid x)\prod_{x'\ne x}\sum_{y'\in S_Y}
\kappa_Y(y'\mid x')=\kappa_Y(y\mid x)\). By
\cref{thm:p16-causal-factorization} the observational law is
\(p_Z(z)\kappa_X(x\mid z)\kappa_Y(y\mid x)\).

If \(p_{ZX}(z,x)>0\), then \(p_Z(z)>0\) and \(p_X(x)>0\), and the product
is
\[
 p_Z(z)\cdot\frac{p_{ZX}(z,x)}{p_Z(z)}\cdot\frac{p_{XY}(x,y)}{p_X(x)}
 =\frac{p_{ZX}(z,x)\,p_{XY}(x,y)}{p_X(x)}=p(z,x,y),
\]
the last equality being the hypothesis divided by \(p_X(x)>0\). If
\(p_{ZX}(z,x)=0\), then \(p(z,x,y)=0\); either \(p_Z(z)=0\), or
\(p_Z(z)>0\) and \(\kappa_X(x\mid z)=0\), and in both cases the product
is \(0\). So the observational law of the constructed model is \(p\).
\end{proof}
\end{annexitem}
